\documentstyle[% 


righttag% 


,11pt% 


,amssymb% 


,russian%               remove in Eng. 


%,izvru%                 Set izven% in Eng. 


,izvru-ks%             for brief communications 


]{amsart} 


 


%  


\newcommand{\down}[2]{\underset{#1}{#2}{}} 


\newcommand{\const}{\operatorname{const}} 


\newcommand{\diag}{\operatorname{diag}} 


\newcommand{\tts}[1]{} 


 


\theoremstyle{plain} 


\theorembodyfont{\it} 


\newtheorem{thmnonum}{\hskip\parindent Теорема} 


\renewcommand{\thethmnonum}{} 


 


%\hoffset=-15mm%                  For Eng. 


 


\setcounter{page}{79} 


\god{2000} 


\nomer{9~(460)} %                        Remove in Eng. 


%\nomer{Vol.\,44, No.\,9}%               Set in Eng. 


%\str{\pageref{begin}--\pageref{end}}%  Set in Eng. 


\udk{514.7} 


 


\author{\it С.Г.\,Лейко, А.В.\,Винник} 


% NB:   ^^ remove in Eng. 


\title{Поворотно-конформные преобразования\\ в  


плоскости Лобачевского} 


 


\date{} 


% \pagestyle{myheadings}%         Set in Eng. 


\pagestyle{plain} %              Remove in Eng. 


\begin{document} 


% \thispagestyle{plain}%          Set in Eng. 


\maketitle 


\label{begin} 


 


В двумерном римановом пространстве $(M^2,g)$ гауссовой кривизны $K\ne0$ 


кривые, геодезическая кривизна которых удовлетворяет условию $k_g=c{}K$, 


$c=\const$, являются изопериметрическими экстремалями поворота (ИЭП). 


Преобразования пространства или его части, переводящие геодезические 


кривые в ИЭП, названы поворотными. Аналогично, если локальное 


инфинитезимальное преобразование $\wt x^h=x^h+\varepsilon X^h(x)$ 


переводит каждую 


геодезическую в кривую, которая в главном (относительно параметра


$\varepsilon$) 


является ИЭП, то оно называется инфинитезимальным поворотным 


преобразованием [1]--[3]. Геодезические преобразования являются 


тривиальными поворотными преобразованиями. 


 


Пусть векторное поле $X$ определяет инфинитезимальное конформное 


преобразование, т.\,е.


\begin{equation} 


L_X g_{ij}=\nabla_i X_j+\nabla_j X_i=2\varphi g_{ij}. 


\end{equation} 


Это преобразование будет поворотным только в случае, когда 


соответствующая функция конформности $\varphi(x)$ порождает специальное 


конциркулярное ковекторное поле $\varphi_i=\partial_i\varphi$:


\begin{equation} 


\nabla_j\varphi_i=\varphi_i\partial_j\ln|K|+(C-\varphi)K g_{ij},\quad 


C\text{ --- }\const. 


\end{equation} 


При $\varphi=C$ преобразование будет гомотетическим и, следовательно, 


геодезическим (т.\,е. тривиальным поворотным преобразованием). Если


$\varphi$ отлична от константы, пространство $(M^2,g)$ является локально 


изометричным поверхности вращения $x=r\cos v$, $y=r\sin v$, $z=f(r)$, 


$$ 


(g_{ij})=\diag\bigg(1+\bigg(\frac{df}{dr} \bigg)^2,\ r^2 \bigg). 


$$ 


Интегрирование (1), (2) на поверхности вращения дает 


$$ 


C=0,\quad \varphi=\frac{c_1} 


{\sqrt{1+\left(\frac{df}{dr}\right)^2}},\quad 


X^h=\bigg(\frac{c_1r}{\sqrt{1+\left(\frac{df}{dr}\right)^2}}, 


\ c_2\bigg),\quad c_1,\,c_2\text{ --- константы}. 


$$ 


В частности, на псевдосфере радиуса $R$, где локально реализуется 


плоскость Лобачевского 


$(g_{ij})=\begin{pmatrix} 


\frac{R^2}{r^2} & 0 \\ 0 & r^2 


\end{pmatrix}$, 


$$ 


\frac{df}{dr}=\frac{\sqrt{R^2-r^2}}{r},\quad 


X^h=\bigg(\frac{c_1r^2}{R},\ c_2 \bigg),\quad 


K=-\frac{1}{R^2}. 


$$ 


Векторные поля $\down{1}{X}^h=\big(\frac{r^2}{R},0\big)$, 


$\down{2}{X}^h=(0,1)$ являются базисом операторов


двучленной группы локальных 


преобразований псевдосферы: $\wt v=v+\tau$, 


$\wt r=\frac{R{}r}{R-rt}$. 


Эти преобразования будут поворотными. Действительно, как было 


показано в [1], функция конформности 


$e^{\psi(\wt r)}:dl^2=e^{2\psi(\wt r)}d\wt l^2$ 


в общих по 


отображению координатах $\wt r$, $\wt v$ должна иметь вид 


$$ 


e^{\psi(\wt r)}=-\frac{\sqrt{1+\left(\frac{df}{d\wt r}\right)^2}} 


{A\left(B+\sqrt{1+\left(\frac{df}{d\wt r}\right)^2}\right)}, 


\ \; A\ne0,\ \; B\text{ --- }\const,\quad 


\text{или}\quad 


e^{\psi(\wt r)}=B\sqrt{1+\left(\tfrac{df}{d\wt r}\right)^2}, 


\ \; B\text{ --- }\const. 


$$ 


Непосредственная проверка показывает, что здесь имеет место 


первый случай при $A=-1$, $B=t$. 


 


Аналогично, интегрирование уравнений (1), (2) в полуплоскости 


Пуанкаре: $x^1{=}x$, $x^2{=}y{>}0$, 


$g_{ij}=\frac{1}{y^2}\delta_{ij}$, дает $C=0$, 


\begin{gather*} 


\varphi=\frac{1}{y}


\bigg[\frac{a_0}{2}x^2+a_1x+c_1\bigg]+\frac{a_0}{2}y,\\ 


% 


X^h=\bigg(y(a_0x+a_1)+\frac{k}{2}(x^2-y^2)+ 


k_1x+k_0,\ y\bigg(\frac{a_0}{2}y+k x+k_1\bigg)- 


\frac{a_0}{2}x^2-a_1x-c_1\bigg),\\ 


% 


a_0,\ a_1,\ k_0,\ k_1,\ k,\ c_1\;\text{ --- константы}. 


\end{gather*} 


Bекторные поля


\begin{alignat*}{3} 


\down{1}{X}&= \begin{pmatrix} 


x y \\ \frac{y^2-x^2}{2} 


\end{pmatrix} 


, &\quad \down{2}{X}&=\begin{pmatrix} 


y \\ -x 


\end{pmatrix} 


, &\quad \down{3}{X}&=\begin{pmatrix} 


0 \\ -1 


\end{pmatrix} 


, \\ 


% 


\down{4}{X}&=\begin{pmatrix} 


1 \\ 0 


\end{pmatrix} 


, &\quad \down{5}{X}&=\begin{pmatrix} 


x \\ y 


\end{pmatrix} 


, &\quad \down{6}{X}&=\begin{pmatrix} 


\frac{x^2-y^2}{2} \\ x y 


\end{pmatrix} 


\end{alignat*} 


являются базисом операторов группы поворотно-конформных 


преобразований плоскости Лобачевского в модели Пуанкаре. 


Как известно ([4], с.\,33), группа конформных преобразований Мебиуса 


$\wt z=\frac{az+b}{cz+d}$, $z=x+iy$, где $a$, $b$, $c$, $d$ --- 


комплексные постоянные, имеет 


следующие базисные операторы: 


\begin{alignat*}{3} 


\down{1}{I}&=\begin{pmatrix} 


\frac{1-x^2+y^2}{2} \\ -x y 


\end{pmatrix} 


, &\quad \down{2}{I}&=\begin{pmatrix} 


\frac{-1-x^2+y^2}{2} \\ -x y 


\end{pmatrix} 


, &\quad \down{3}{I}&=\begin{pmatrix} 


x \\y 


\end{pmatrix} 


, \\ 


\down{4}{I}&=\begin{pmatrix} 


-x y \\\frac{1+x^2-y^2}{2} 


\end{pmatrix} 


, &\quad \down{5}{I}&=\begin{pmatrix} 


-x y \\ \frac{-1+x^2-y^2}{2} 


\end{pmatrix} 


, &\quad \down{6}{I}&=\begin{pmatrix} 


y \\ -x 


\end{pmatrix} 


. 


\end{alignat*} 


Нетрудно видеть, что 


\begin{alignat*}{3} 


\down{1}{X}&=-\frac{1}{2}(\down{4}{I}+\down{5}{I}), &\quad 


\down{2}{X}&=\down{6}{I}, &\quad \down{3}{X}&= 


\down{5}{I}-\down{4}{I}, \\ 


% 


\down{4}{X}&=\down{1}{I}-\down{2}{I}, &\quad 


\down{5}{X}&=\down{3}{I}, &\quad \down{6}{X}&= 


-\frac{1}{2}(\down{1}{I}+\down{2}{I}). 


\end{alignat*} 


Таким образом, операторы 


$\down{1}{X}$, $\down{2}{X}$, $\down{3}{X}$, 


$\down{4}{X}$, $\down{5}{X}$, $\down{6}{X}$ 


также образуют базис 


указанной 


группы конформных преобразований Мебиуса в плоскости Лобачевского. 


Отсюда вытекает 


 


\begin{thmnonum} 


Группа поворотно-конформных преобразований плоскости 


Лобачевского совпадает с группой конформных преобразований Мебиуса 


этой плоскости. 


\end{thmnonum} 


 


Отметим, что операторы 


$\down{4}{X}$, $\down{5}{X}$, $\down{6}{X}$ 


определяют движение в плоскости 


Лобачевского (случай вещественных параметров $a$, $b$, $c$, $d$). 


Оставшиеся 


операторы 


$\down{1}{X}$, $\down{2}{X}$, $\down{3}{X}$ 


выделяют группу нетривиальных 


поворотно-конформных преобразований. 


 


\begin{thebibliography}{9} 


 


\bibitem{1} 


Лейко С.Г. {\em Поворотные диффеоморфизмы на поверхностях 


евклидового пространства} // Матем. заметки. -- 1990. 


-- \No\,3. -- С.\,52--57. 


\bibitem{2} 


Лейко С.Г. {\em Инфинитезимальные поворотные преобразования и 


деформации поверхностей евклидового пространства} // Докл. 


РАН. -- 1995. -- Т.\,344. -- \No\,2. -- С.\,162--164. 


\bibitem{3} 


Лейко С.Г. {\em Поворотные преобразования поверхностей} // Матем. 


физика, анализ, геометрия. -- 1998. -- Т.5. 


-- \No\,3/4. -- С.\,203--211. 


\bibitem{4} 


Широков А.П. {\em Неевклидовы пространства}. Учеб. пособие. -- 


Казань: Изд-во Казанск. ун-та, 1997. -- 49 с. 


\bibitem{5} 


Букреев Б.Я. {\em Планиметрия Лобачевского в аналитическом 


изложении}. -- М.--Л.: ГИТТЛ, 1951. -- 127 с. 


 


\end{thebibliography} 


 


\vskip 2cm 


 


\postup{\it Одесский государственный университет}% 


{$\qquad$\it Поступили} 


\postup{}{\it первый вариант $10.03.1999$} 


\postup{}{\it окончательный вариант $07.02.2000$} 


\label{end} 


\end{document} 


